%%%%%%%%%%%%%%%%%%%%%%%%%%%%%%%%%%%%%%%%%
% a0poster Portrait Poster for LiRI/UZH
% LaTeX Template
% Version 1.0 (22/06/13)
% Version 2.0 (08/04/22)
%
% The a0poster class was created by:
% Gerlinde Kettl and Matthias Weiser (tex@kettl.de)
% adapted by Danny McDonald for LiRI/UZH (mcddjx@gmail.com)
%
% License:
% CC BY-NC-SA 3.0 (http://creativecommons.org/licenses/by-nc-sa/3.0/)
%
%%%%%%%%%%%%%%%%%%%%%%%%%%%%%%%%%%%%%%%%%

\documentclass[a0,portrait]{a0poster}

% control margins here
\usepackage{geometry}
 \geometry{
 a0paper,
 left=5cm,
 top=4.5cm,
 right=5cm,
 bottom=4.5cm
 }
\addtolength{\textwidth}{4.5cm} % width of text can be adjusted if margins above are adjusted

\usepackage{multicol} % This is so we can have multiple columns of text side-by-side
\columnsep=3em % This is the amount of white space between the columns in the poster
\columnseprule=0pt % This is the thickness of the black line between the columns in the poster

% UZH colours
\usepackage[svgnames]{xcolor}
\definecolor{red100}{RGB}{204, 3, 1}
\definecolor{uzhblau80}{RGB}{51,83,183}
\definecolor{uzhockerrot100}{RGB}{220, 96, 39}
\definecolor{uzhockerrot80}{RGB}{227, 128, 82}
\definecolor{uzhflaschengruen100}{RGB}{42, 127, 98}
\definecolor{uzhflaschengruen80}{RGB}{86, 157, 133}
\definecolor{about}{RGB}{204,237,212} % 
\definecolor{conclusion}{RGB}{204,212,237} % the conclusion box colour
\usepackage{multirow}

\usepackage{ifthen} % needed to stop horizontal line above 'Conclusions' section
\usepackage{graphicx} % Required for including images
\graphicspath{{figures/}} % Location of the graphics files
\usepackage{mwe,tikz}\usepackage[percent]{overpic} % overlay your photo over the background
\usepackage{booktabs} % Top and bottom rules for table
\usepackage[font=small,labelfont=bf]{caption} % Required for specifying captions to tables and figures
\usepackage{amsfonts, amsmath, amsthm, amssymb} % For math fonts, symbols and environments
\usepackage{wrapfig} % Allows wrapping text around tables and figures
\usepackage{tgheros}
\renewcommand{\sfdefault}{qhv}
\renewcommand{\familydefault}{\sfdefault}

\usepackage[onehalfspacing]{setspace} % remove for single spacing
\usepackage{tcolorbox} % for the conclusions box
\usepackage{blindtext} % you can remove this once you add your content
\usepackage[export]{adjustbox} % allow floating a graphic right
\usepackage{titlesec} % customising section titles
\usepackage{needspace} % prevent break between line and section title
\usepackage{nameref} % package and command to get the name of the current section (for ifthen)
\usepackage[version=4]{mhchem}
\usepackage{tabularx}
\usepackage{booktabs}
\newcolumntype{R}[1]{>{\raggedleft\arraybackslash}p{#1}}

\begin{document}

% define how our section titles will look (with ruled line)
\titleformat{\section}
  {\needspace{1\baselineskip}\sectionrule\huge\bfseries}
  {\color{red100}\thesection.}
  {1em}
  {\color{red100}}

% draw horizontal line before section unless it is conclusions (if you change name of Conclusions, you should
% also change it here too so it is recognised and the line suppressed
\makeatletter
\newcommand{\sectionrule}{%
 \ifthenelse{\equal{\@currentlabelname}{Conclusions}}
 % use the below line instead of the above if conclusions is a section*
 % \ifthenelse{\equal{\@currentlabelname}{}}
  {}
  {\vspace*{-\baselineskip}
  %  \vrule height 1pt depth 1pt width \linewidth\vskip0.4pt
   \bigskip}%
}
\makeatother


%----------------------------------------------------------------------------------------
%	POSTER HEADER 
%----------------------------------------------------------------------------------------
%\title{Creating a \LaTeX{} poster template matching UZH specifications for LiRI presentations}
\title{The Role of PCI-PMI Infrastructure in Reaching Europe's Energy Policy Targets}

% the top logo (in english!) and unit title
% \noindent

\noindent
\begin{minipage}[b][2cm][t]{0.753\linewidth} 
  % First Row
  \vspace{1em}
  \begin{minipage}{0.15\textwidth}
    \includegraphics[height=3.5cm]{tub-logo.pdf}
  \end{minipage}
  \begin{minipage}{0.2\textwidth}
    \includegraphics[height=4.5cm]{ensys-logo-en.pdf}
  \end{minipage}
  \vspace{2.3em}

  % Second Row
  \begin{minipage}{0.9\textwidth}
    \makeatletter
    \raggedright{\fontsize{75pt}{65pt}\selectfont\color{red100}\textbf{{\@title}}\par}
    \makeatother
    \color{Black}
    \vspace{0.3cm}
    \hspace{0.005em}
    resilient-project.github.io $\cdot$ pypsa.org $\cdot$ github.com/PyPSA/pypsa-eur $\cdot$ github.com/bobbyxng
  \end{minipage}
\end{minipage}%
%
\begin{minipage}[b][2cm][t]{0.3\linewidth} % Second minipage
    \begin{minipage}[t]{0.3\linewidth}
        \includegraphics[height=13cm]{resilient_bmwk_logo.png}
    \end{minipage}
    \hfill
    \begin{minipage}[t]{0.7\linewidth}
        \includegraphics[height=13cm]{qr.pdf}
    \end{minipage}%     
\end{minipage}

\vspace{15em}



%----------------------------------------------------------------------------------------
%	POSTER BODY
%----------------------------------------------------------------------------------------

\begin{multicols}{3}
\raggedcolumns

%----------------------------------------------------------------------------------------
%	PYPSA
%----------------------------------------------------------------------------------------

\noindent \textcolor{red100}{\huge \textbf{PCI-PMI projects}}
\\
\textit{\ce{CO2} and \ce{H2} infrastructure}
\\
\noindent
\begin{center}
    \includegraphics[width=\linewidth]{map_adm_pcipmi.pdf}
    \captionof{figure}{Map of the regional scope including clustered onshore (grey) and offshore regions (blue), as well as PCI-PMI \ce{CO2} and \ce{H2} pipelines, storage and sequestration sites. Depleted offshore oil and gas fields (red) provide additional \ce{CO2} sequestration potential \cite{hofmannH2CO2Network2025}.}
    \label{fig:regional_scope_map}
\end{center}
\vspace{2em}

\begin{center}
    \captionof{table}{Regret matrix setup: Long-term and short-term scenarios.}
    \label{fig:regret_matrix_setup}
    \scriptsize
    \begin{tabularx}{\linewidth}{R{7.5cm}>{\centering\arraybackslash}X>{\centering\arraybackslash}X>{\centering\arraybackslash}X}
        \toprule
        \textbf{Short-term} & \textbf{Reduced targets} & \textbf{Delayed pipelines} & \textbf{No pipelines} \\
        \midrule
        \textbf{Long-term scenarios} & & & \\
        Decentral Islands (\textbf{DI}) & $\blacksquare$ & -- & -- \\
        PCI-PMI (\textbf{PCI}) & $\blacksquare$ & $\blacksquare$ & $\blacksquare$ \\
        PCI-PMI nat. (\textbf{PCI-n}) & $\blacksquare$ & $\blacksquare$ & $\blacksquare$\\
        PCI-PMI internat. (\textbf{PCI-in}) & $\blacksquare$ & $\blacksquare$ & $\blacksquare$ \\
        Central Planning (\textbf{CP}) & $\blacksquare$ & $\blacksquare$ & $\blacksquare$ \\
        \midrule
        \textbf{Targets} & & & \\
        GHG emission reduction &  $\blacksquare$ &  $\blacksquare$ &  $\blacksquare$ \\
        \ce{CO2} sequestration &  -- &  $\blacksquare$ &  $\blacksquare$ \\
        Electrolytic \ce{H2} production &  -- &  $\blacksquare$ &  $\blacksquare$ \\
        \ce{H2} electrolysers &  -- &  $\blacksquare$ &  $\blacksquare$ \\
        \midrule
        \textbf{\ce{CO2} + \ce{H2} infrastructure} & & & \\
        \ce{CO2} sequestration sites & $\blacksquare$ &  $\blacksquare$ &  $\blacksquare$ \\
        \ce{CO2} pipelines to seq. site & $\blacksquare$ &  $\blacksquare$ &  $\blacksquare$ \\
        \ce{CO2} pipelines & $\blacksquare$ &  $\square$ &  -- \\
        \ce{H2} pipelines & $\blacksquare$ &  $\square$ &  -- \\
        \bottomrule
    \end{tabularx}
\end{center}

\columnbreak

\noindent \textcolor{red100}{\huge \textbf{Results \& Discussion}}
\\
\textit{\ce{CO2} and \ce{H2} infrastructure}
\\

\begin{center}
    \includegraphics[width=\linewidth]{base_s_adm___2050-balance_map_co2_stored.jpg}
    \captionof{figure}{\textit{PCI-in long-term} scenario (2050) --- \ce{CO2} balances and distribution}
    \label{fig:PCI-in_lt_2050_co2}
\end{center}
\vspace{2em}
\begin{center}
    \includegraphics[width=\linewidth]{base_s_adm___2050-balance_map_H2.jpg}
    \captionof{figure}{\textit{PCI-in long-term} scenario (2050) --- \ce{H2} balances and distribution}
    \label{fig:PCI-in_lt_2050_h2}
\end{center}


% \noindent \textcolor{red100}{\huge \textbf{PyPSA}}
% \\
% \textit{Python for Power System Analysis}
% \\

% \noindent
% PyPSA is an open-source toolbox for \textbf{simulating and optimising modern power and energy 
% systems} that include features such as conventional generators with unit commitment, 
% variable wind and solar generation, storage units, coupling to other energy sectors, 
% and mixed alternating and direct current networks. \\

% \noindent \textcolor{black}{\large \textbf{Functionality}}

% \noindent PyPSA can calculate:

% \begin{itemize}
%     \item static power flow (using both the full non-linear network equations and the linearised network equations)
%     \item linear optimal power flow (least-cost optimisation of power plant and storage dispatch within network constraints, using the linear network equations, over several snapshots)
%     \item security-constrained linear optimal power flow
%     \item total electricity/energy system least-cost investment optimisation
% \end{itemize}

\columnbreak
%----------------------------------------------------------------------------------------
%	SECTIONS
%----------------------------------------------------------------------------------------
\noindent \textcolor{red100}{\huge \textbf{Conclusion \& Key Findings}}
\\
\textit{TODO FIll}
\\


\begin{center}
    \includegraphics[width=\linewidth]{regret_matrix.pdf}
    \captionof{figure}{Regret matrix. Calculating regret terms by subtracting system costs of long-term scenarios (columns) from short-term scenarios (rows). Positive values reflect higher costs in the short-term scenarios compared to the long-term ones.}
    \label{fig:regret_matrix}
\end{center}
\vspace{2em}

\noindent
PyPSA-Eur is a spatially and temporally highly resolved \textbf{linear optimisation 
model that covers the European continent}. The model is suitable both for operational studies and 
generation and transmission expansion planning studies. The continental scope and 
highly resolved spatial scale enables a proper description of the long-range smoothing
effects for renewable power generation and their varying resource availability.

\noindent

\vspace{2em}
\begin{tcolorbox}[width=0.95\linewidth,colback={about},frame empty,boxsep=1cm]
\section*{RESILIENT Project}
    Resilient Energy System Infrastructure Layouts for Industry, E-Fuels and Network Transitions 
\end{tcolorbox}   

\vspace{2em}
\begin{tcolorbox}[width=0.95\linewidth,colback={conclusion},frame empty,boxsep=1cm]
\section*{Open Source}
    All tools are open-source and available on GitHub. The code is licensed under the MIT
    license, which allows for free use and modification of the code. 
    PyPSA and PyPSA-Eur are actively developed and maintained by the Department of Digital
    Transformation in Energy Systems at the TU Berlin and the global \textbf{PyPSA community}.
\end{tcolorbox}    

%---------------------------------------------------------------------------------------
%	REFERENCES
%---------------------------------------------------------------------------------------

\singlespacing
\small
% \nocite{*} % Print all references regardless of whether they were cited in the poster or not
\bibliographystyle{plain} % Plain referencing style
\bibliography{references}

%---------------------------------------------------------------------------------------
\end{multicols}

% Add keywords and author in footer minipage
\vfill % pushes footer to bottom if space allows
\begin{center}
  \colorbox{red100}{\parbox{1\linewidth}{
    \centering
    \color{white}
    \textbf{Contact:} Bobby Xiong (xiong@tu-berlin.de) \quad
    \textbf{Keywords:} \quad energy system modelling $\cdot$ policy targets $\cdot$ infrastructure $\cdot$ resilience $\cdot$ hydrogen $\cdot$ carbon $\cdot$ Europe 
  }}
\end{center}

\end{document}